\section{Conclusion}\label{sec:conclusion}

Clustering failed programs by the tests they fail is an appealing way to give instructors macro-feedback, and its IF--THEN rules read like teaching advice. Evaluated against verified minimal repairs of the students' own programs and against controlled single-mechanism mutants, the test-case OAV turned out to mix mechanisms systematically: it identifies the dominant output-level failures and merges much of the rest (RQ1). Label-free evidence about how the output deviates from the oracle recovers more, most clearly where mechanisms are balanced (RQ2), and it supports rules that carry across assignments (RQ3); general-purpose code embeddings, in contrast, cluster programs by origin (RQ4). Because every label in our benchmarks is derived from executed evidence and can be re-derived, the benchmarks can be reused to evaluate new representations, LLM-based diagnosers or active tests, and the instructor tool shows how the findings translate into feedback that abstains when the evidence is weak.

Three steps would strengthen these conclusions: an independent audit of the labels by two human raters, replication on other courses, languages and grading policies, and a classroom study of whether mechanism hypotheses change what instructors re-teach. Adaptive tests that break the remaining signature collisions before a hypothesis reaches the instructor are a natural next method. Claims about students' beliefs need studies with students; the tools released here are meant to make such studies cheaper to design.
